\documentclass[17pt]{article}

\usepackage{cmap}	% поиск по документу
\usepackage{hyperref}	% поддержка гиперссылок
\usepackage{geometry}	% настройка геометрии документа
\usepackage{multicol}   % настройка шаблона расположения текста
\usepackage{extsizes}	% расширенная настройка размеров шрифтов
\usepackage{indentfirst}    % отступ для первого абзаца
\usepackage{tabularx}	% расширенные таблицы
\usepackage{makecell}   % настройка клеток таблиц
\usepackage{caption}    % настройка подписей к вставкам
\usepackage[table]{xcolor}  % покраска таблиц
\usepackage{graphicx}	% вставка изображений
\usepackage{wrapfig, booktabs}	% оформление таблиц
\usepackage{graphicx, caption, subcaption}
\usepackage{minted}    % вставка кода

\usepackage{color}	% настройка цвета текста - \textcolor
\usepackage{soul}	% выделение текста - \hl
\usepackage[T2A]{fontenc}	% поддержка кириллических символов
\usepackage[utf8]{inputenc} % поддержка кириллических символов
\usepackage[english,russian]{babel}	% поддержка кириллических символов
\usepackage{csquotes, biblatex}   % вывод библиографии

\usepackage{mathcmd}	% математические команды
\usepackage{amsfonts}	% математические символы 1
\usepackage{amssymb}	% математические символы 2
\usepackage{amsmath}	% математические символы 3
\usepackage{amsthm}	% работа с теоремами

% Настройка окружения теорем 1 (RU)
\newtheorem{theorem}{Теорема}[section]
\newtheorem{corollary}{Следствие}[theorem]
\newtheorem{lemma}[theorem]{Лемма}
\newtheorem{example}{Пример}[section]
\newtheorem*{remark}{Утверждение}
\addto\captionsrussian{\renewcommand\proofname{Док-во:}}

% Настройка окружения теорем 2 (EN)
%\newtheorem{theorem}{TH}[section]
%\newtheorem{corollary}{CO}[theorem]
%\newtheorem{lemma}[theorem]{LM}
%\newtheorem{example}{EX}[section]
%\addto\captionsrussian{\renewcommand\proofname{PF:}}

\renewcommand{\labelitemii}{$\circ$}
\renewcommand{\labelitemiii}{-}

\addbibresource{biblio.bib}
\graphicspath{ {./images/part2} }

% Настройка оформления вставок кода
\setminted{xleftmargin=\parindent, fontsize=\footnotesize, linenos, tabsize=2, breaklines}
\usemintedstyle{emacs} 

% Настройка отступов
\geometry{
    a4paper,
    % total={200mm,287mm},
    left=10mm,
    right=10mm,
    top=10mm,
    bottom=20mm
}

% Заголовок
\title{\large{Методы моделирования и анализа стохастических систем\\ Лабораторная работа № 3, Часть 2}}
\author{\large{Костарев Иван, МЕНМ-150201}}
\date{весенний семестр 2025/2026}


\begin{document}

\maketitle
\tableofcontents
\newpage



\section*{Исследование разброса случайных состояний системы вблизи устойчивого равновесия}

\section{Статистический анализ результатов прямого численного моделирования}

(а) Определим процедуры обновления состояния \texttt{step\_stochastic\_em}, \texttt{step\_stochastic\_rk4}, задающие временной шаг стохастической системы, и запишем стохастические системы FHN и SN. В работе будем использовать стохастический метод Рунге-Кутты 4-го порядка для обновления состояния (процедура \texttt{step\_stochastic\_rk4}).

FHN с аддитивным шумом в 1-м уравнении:  
$$
\begin{cases}
    \dot{x} = \frac{1}{\delta} (x - \frac{x^3}{3} - y) + \varepsilon \cdot \xi_1 \\
    \dot{y} = x + a
\end{cases}
$$

SN с аддитивным шумом в обоих уравнениях:  
$$
\begin{cases}
    \dot{x} = y - x + \varepsilon \cdot \xi_1 \\
    \dot{y} = -a x + b y - x^2 y + \varepsilon \cdot \xi_2
\end{cases}
$$

(б) Для устойчивых равновесий построим облака случайных траекторий для разных значений интенсивности шума $\varepsilon$.

На основе симуляции длительностью $T = 100$ с шагом по времени $h = 0.01$, получены облака траекторий системы FHN вблизи устойчивого равновесия $(\overline{x}, \overline{y}) = (-a, \frac{a^3}{3} - a)$ при $a > 1$ (Рис. \ref{fig:stochastic_clouds_FHN}).

\begin{figure}[H]
    \centering
    \includegraphics[width=1\textwidth]{stochastic_clouds_FHN.png}
    \caption{Облака траекторий системы FHN}
    \label{fig:stochastic_clouds_FHN}
\end{figure}

На основе симуляции длительностью $T = 100$ с шагом по времени $h = 0.01$, получены облака траекторий системы SN вблизи устойчивого равновесия $(\overline{x}, \overline{y}) = (\sqrt{b - a}, \sqrt{b - a})$ при $a < 1$ (Рис. \ref{fig:stochastic_clouds_SN}).

\begin{figure}[H]
    \centering
    \includegraphics[width=1\textwidth]{stochastic_clouds_SN.png}
    \caption{Облака траекторий системы SN}
    \label{fig:stochastic_clouds_SN}
\end{figure}

(в) Запишем матрицу ковариации для облака при фиксированном $\varepsilon$, найдем ее собственные числа $\overline{\lambda_1}(a), \overline{\lambda_2}(a)$.

$$
\operatorname{cov}(x, y) =
\begin{pmatrix}
    \frac{1}{n} \sum_{i=1}^n (x_i - \bar{x})^2 & \frac{1}{n} \sum_{i=1}^n (x_i - \bar{x})(y_i - \bar{y}) \\
    \frac{1}{n} \sum_{i=1}^n (x_i - \bar{x})(y_i - \bar{y}) & \frac{1}{n} \sum_{i=1}^n (y_i - \bar{y})^2
\end{pmatrix}
$$

Построим графики зависимости собственных чисел матрицы ковариации облака траекторий от параметра $a$ ($\overline{\lambda_1}(a), \overline{\lambda_2}(a)$) в зоне устойчивых равновесий $(a_1, a_2)$.

Будем использовать от $1$ до $100$ симуляций и $10^3$ шагов по времени в каждой из них, таким образом получим графики для выборок размеров $N = 1 \cdot 10^3$ и $N = 100 \cdot 10^3$ при значении $\varepsilon = 0.05$.

В модели FHN при значении параметра $a > 1$ равновесие $(x, y) = (0, 0)$ устойчиво. На основе симуляции длительностью $T = 10$ с шагом по времени $h = 0.01$ построим облако траекторий, вычислим матрицу ковариации и изобразим эмпирическую зависимость $\overline{\lambda_1}(a), \overline{\lambda_2}(a)$ (Рис. \ref{fig:eigenvalues_empirical_FHN}).

\begin{figure}[H]
    \centering
    \includegraphics[width=1\textwidth]{eigenvalues_empirical_FHN.png}
    \caption{Собственные значения матрицы ковариации для модели FHN}
    \label{fig:eigenvalues_empirical_FHN}
\end{figure}

В модели SN при значении параметра $a < 1$ устойчивыми являются равновесия:

$$
\begin{cases}
    (\overline{x}_2, \overline{y}_2) = (\sqrt{b - a}, \sqrt{b - a}) \\
    (\overline{x}_3, \overline{y}_3) = (-\sqrt{b - a}, -\sqrt{b - a})
\end{cases}
$$

Вблизи них система ведет себя одинаково с точностью до симметрии. Рассмотрим одно из них, например $(\overline{x}_2, \overline{y}_2)$. На основе симуляции длительностью $T = 10$ с шагом по времени $h = 0.01$ построим облако траекторий, вычислим матрицу ковариации и изобразим эмпирическую зависимость $\overline{\lambda_1}(a), \overline{\lambda_2}(a)$ (Рис. \ref{fig:eigenvalues_empirical_SN}).

\begin{figure}[H]
    \centering
    \includegraphics[width=1\textwidth]{eigenvalues_empirical_SN.png}
    \caption{Собственные значения матрицы ковариации для модели SN}
    \label{fig:eigenvalues_empirical_SN}
\end{figure}



\newpage
\section{Анализ с помощью метода стохастической чувствительности}
(а) Запшем решение уравнения Ляпунова в общем виде.

$$
F W + W F^T + S S^T = 0
$$

С учетом изначальной формы СДУ:

$$
\begin{cases}
    \dot{x} = f(x, y) + \varepsilon \sigma_1(x, y) \xi_1 \\
    \dot{y} = g(x, y) + \varepsilon \sigma_2(x, y) \xi_2
\end{cases}
$$

Имеем
$$
S = 
\begin{pmatrix}
    \sigma_1(x, y) & 0 \\
    0 & \sigma_2(x, y)
\end{pmatrix}
\implies
SS^T = 
\begin{pmatrix}
    \sigma_1^2(x, y) & 0 \\
    0 & \sigma_2^2(x, y)
\end{pmatrix}
= 
\begin{pmatrix}
    s_1 & 0 \\
    0 & s_2
\end{pmatrix}
$$

$$
F = 
\begin{pmatrix}
    f'_x(x, y) & f'_y(x, y) \\
    g'_x(x, y) & g'_y(x, y)
\end{pmatrix}
=
\begin{pmatrix}
    a & b \\
    c & d
\end{pmatrix}
$$

$$
W = 
\begin{pmatrix}
    w_1 & w_2 \\
    w_2 & w_3
\end{pmatrix}
$$

Тогда уравнение (3) имеет вид:

$$
\begin{pmatrix}
    a & b \\
    c & d
\end{pmatrix}
\begin{pmatrix}
    w_1 & w_2 \\
    w_2 & w_3
\end{pmatrix}
+
\begin{pmatrix}
    w_1 & w_2 \\
    w_2 & w_3
\end{pmatrix}
\begin{pmatrix}
    a & c \\
    b & d
\end{pmatrix}
+
\begin{pmatrix}
    s_1 & 0 \\
    0 & s_2
\end{pmatrix}
= 0
$$

$$
\begin{pmatrix}
    a w_1 + b w_2 & a w_2 + b w_3 \\
    c w_1 + d w_2 & c w_2 + d w_3
\end{pmatrix}
+
\begin{pmatrix}
    a w_1 + b w_2 & c w_1 + d w_2 \\
    a w_2 + b w_3 & c w_2 + d w_3
\end{pmatrix}
=
- \begin{pmatrix}
    s_1 & 0 \\
    0 & s_2
\end{pmatrix}
$$

$$
\begin{pmatrix}
    2a w_1 + 2b w_2 & cw_1 + (a+d)w_2 + bw_3 \\
    cw_1 + (a+d)w_2 + bw_3 & 2c w_2 + 2d w_3
\end{pmatrix}
=
- \begin{pmatrix}
    s_1 & 0 \\
    0 & s_2
\end{pmatrix}
$$

В силу симметрии матриц в левой и правой частях, получаем систему из 3-х уравнений:
$$
\begin{cases}
   2a w_1 + 2b w_2 = -s_1 \\
   c w_1 + (a+d) w_2 + b w_3 = 0 \\
   2c w_2 + 2d w_3 = -s_2
\end{cases}
\implies
\begin{cases}
    w_1 = -\frac{s_1 + 2b w_2}{2a} \\
    w_3 = -\frac{s_2 + 2c w_2}{2d} \\
    -c \cdot \frac{s_1 + 2b w_2}{2a} + (a + d) w_2 - b \cdot \frac{s_2 + 2c w_2}{2d} = 0
\end{cases}
$$

Решение относительно $w_1, w_2, w_3$:

$$
\begin{cases}
    w_2 = \frac{
        \frac{c s_1}{2a} + \frac{b s_2}{2c}
    }{
        a + d - \frac{b c}{a} - \frac{b c}{d}
    } \\
    w_1 = -\frac{s_1 + 2b w_2}{2a} \\
    w_3 = -\frac{s_2 + 2c w_2}{2d}
\end{cases}
$$

Получено явное решение, но для удобства программного решения можно представить полученную систему уравнений в матричной форме:

$$
\begin{pmatrix}
    2a & 2b & 0 \\
    c & a+d & b \\
    0 & 2c & 2d
\end{pmatrix}
\begin{pmatrix}
    w_1 \\ w_2 \\ w_3
\end{pmatrix}
=
\begin{pmatrix}
    -s_1 \\ 0 \\ -s_2
\end{pmatrix}
$$

Запишем матрицу стохастической чувствительности $W$ для FHN, где

$$
Q = SS^T = 
\begin{pmatrix}
    \sigma_1^2(x, y) & 0 \\
    0 & \sigma_2^2(x, y)
\end{pmatrix}
=
\begin{pmatrix}
    1 & 0 \\
    0 & 0
\end{pmatrix}
$$

$$
F = 
\begin{pmatrix}
    f'_x(x, y) & f'_y(x, y) \\
    g'_x(x, y) & g'_y(x, y)
\end{pmatrix}
=
\begin{pmatrix}
    \frac{1}{\delta} (1 - x^2) & \frac{1}{\delta} \\
    1 & 0
\end{pmatrix}
$$

Запишем матрицу стохастической чувствительности $W$ для SN, где

$$
Q = SS^T = 
\begin{pmatrix}
    \sigma_1^2(x, y) & 0 \\
    0 & \sigma_2^2(x, y)
\end{pmatrix}
=
\begin{pmatrix}
    1 & 0 \\
    0 & 1
\end{pmatrix}
$$

$$
F = 
\begin{pmatrix}
    f'_x(x, y) & f'_y(x, y) \\
    g'_x(x, y) & g'_y(x, y)
\end{pmatrix}
=
\begin{pmatrix}
    -1 & 1 \\
    -a - 2xy & b - x^2
\end{pmatrix}
$$

% \begin{minted}{python}
% a, b = 0.5, 2
% noise_mask = np.array([1, 1])
% equilibrium_expr = np.array([np.sqrt(b - a), np.sqrt(b - a)])
% W = stochastic_sensitivity_matrix_2d(
%     model_func=model_SN,
%     params=dict(a=a, b=b),
%     equilibrium=equilibrium_expr,
%     noise_mask=noise_mask,
% )
% W
% \end{minted}

(б) Найдем собственные числа $\lambda_1, \lambda_2$ матрицы стохастической чувствительности $W$.

Построим зависимость $\varepsilon \cdot \lambda_1(a), \varepsilon \cdot \lambda_2(a)$ для значения $\varepsilon$ из пункта 1 (б) для системы FHN (Рис. \ref{fig:eigenvalues_theoretical_FHN}).

\begin{figure}[H]
    \centering
    \includegraphics[width=1\textwidth]{eigenvalues_theoretical_FHN.png}
    \caption{Собственные значения матрицы стохастической чувствительности для модели FHN}
    \label{fig:eigenvalues_theoretical_FHN}
\end{figure}

Построим зависимость $\varepsilon \cdot \lambda_1(a), \varepsilon \cdot \lambda_2(a)$ для $\varepsilon$ из пункта 1 (б) для системы SN (Рис. \ref{fig:eigenvalues_theoretical_SN}).

\begin{figure}[H]
    \centering
    \includegraphics[width=1\textwidth]{eigenvalues_theoretical_SN.png}
    \caption{Собственные значения матрицы стохастической чувствительности для модели SN}
    \label{fig:eigenvalues_theoretical_SN}
\end{figure}

Изобразим эмпирическую и теоретическую зависимости на одном графике и сравним их (Рис. \ref{fig:eigenvalues_FHN}, \ref{fig:eigenvalues_SN}).
Эмпирический и теоретический графики схожи по крупномасштабной структуре и по характеру изменения.
Видно, что при увеличении размера выборки эмпирический график становится более гладким и приближается к теоретическому.
Однако, для обеих моделей при приближении к границе интервала стабильности равновесных точек (в правых частях графиков) разница между значениями эмпирического и теоретического графиков растет.

\begin{figure}[H]
    \centering
    \includegraphics[width=1\textwidth]{eigenvalues_FHN.png}
    \caption{Сравнение собственных значений матрицы ковариации и матрицы стохастической чувствительности для модели FHN}
    \label{fig:eigenvalues_FHN}
\end{figure}

\begin{figure}[H]
    \centering
    \includegraphics[width=1\textwidth]{eigenvalues_SN.png}
    \caption{Сравнение собственных значений матрицы ковариации и матрицы стохастической чувствительности для модели SN}
    \label{fig:eigenvalues_SN}
\end{figure}



\newpage
\section{Доверительные эллипсы}
(в) Для выбранных нескольких значений параметра $a$ построим доверительные эллипсы (возьмем облако из пункта 1 (б)).

В двумерном случае матрица стохастической чувствительности $W$ характеризует дисперсию решений стохастической системы вокруг устойчивого равновесия $\overline{x} = (\overline{x_1}, \overline{x_2})$ в форме доверительного эллипса:

$$
\frac{z_1^2}{\lambda_1} + \frac{z_2^2}{\lambda_2} = 2k^2 \varepsilon^2
$$

Пусть $\lambda_1, \lambda_2$ -- собственные числа матрицы $W$, а $V$ -- матрица собственных векторов $v_1 = (v_{11}, v_{21})^T$, $v_2 = (v_{12}, v_{22})^T$, записанных по столбцам:

$$
V = \begin{pmatrix} v_{11} & v_{12} \\ v_{21} & v_{22} \end{pmatrix}
$$

Координаты $x = (x_1, x_2)$ доверительного эллипса в базисе собственных векторов могут быть записаны так:

$$
x = \overline{x} + V \cdot \begin{pmatrix} z_1 \\ z_2 \end{pmatrix},
$$
где $z_1(\varphi) = \sqrt{2\lambda_1} \cdot \varepsilon \cdot k \cdot \cos\varphi$, $z_2(\varphi) = \sqrt{2\lambda_2} \cdot \varepsilon \cdot k \cdot \sin\varphi$ ($0 \leq \varphi \leq 2\pi$).

Коэффициент $k$ связан с доверительной вероятностью $P$ попадания решений стохастической системы в доверительный эллипс:

$$
k = \sqrt{-\ln(1 - P)}
$$

Построим доверительные эллипсы для системы FHN. Рассмотрим несколько значений параметра $a$ внутри интервала, где рассматриваемое равновесие сохраняет устойчивость (Рис. \ref{fig:confidence_ellipse_FHN}).

\begin{figure}[H]
    \centering
    \includegraphics[width=1\textwidth]{confidence_ellipse_FHN.png}
    \caption{Доверительные эллипсы для траекторий модели FHN.}
    \label{fig:confidence_ellipse_FHN}
\end{figure}

Построим доверительные эллипсы для системы SN. Рассмотрим несколько значений параметра $a$ внутри интервала, где рассматриваемое равновесие сохраняет устойчивость (Рис. \ref{fig:confidence_ellipse_SN}).

\begin{figure}[H]
    \centering
    \includegraphics[width=1\textwidth]{confidence_ellipse_SN.png}
    \caption{Доверительные эллипсы для траекторий модели SN.}
    \label{fig:confidence_ellipse_SN}
\end{figure}


% # TODO:
% - Удостовериться, что при симуляции облака траккторий, все траектории стартуют строго в равновесии
% - Построить графики собственных чисел матрицы стохастической чувствительности и матрицы ковариации (теоретический и эмпирический) рядом на одном графике и сравнить

\end{document}
